\documentclass[12pt]{article}
\usepackage{amsmath}
\usepackage{graphicx,psfrag,epsf}
\usepackage{enumerate}
\usepackage{natbib}
\usepackage{url} % not crucial - just used below for the URL


%user supplied additions
\usepackage{amsfonts}
\usepackage{amsthm}
\usepackage{cleveref}
\theoremstyle{plain}
\newtheorem{theorem}{Theorem}[section]
\newtheorem*{theorem*}{Theorem}

%\pdfminorversion=4

% DON'T change margins - should be 1 inch all around.
\addtolength{\oddsidemargin}{-.5in}%
\addtolength{\evensidemargin}{-.5in}%
\addtolength{\textwidth}{1in}%
\addtolength{\textheight}{1.3in}%
\addtolength{\topmargin}{-.8in}%


\begin{document}

%\bibliographystyle{natbib}

\def\spacingset#1{\renewcommand{\baselinestretch}%
{#1}\small\normalsize} \spacingset{1}


%%%%%%%%%%%%%%%%%%%%%%%%%%%%%%%%%%%%%%%%%%%%%%%%%%%%%%%%%%%%%%%%%%%%%%%%%%%%%%

\title{\bf A new tail bound for the sum of bounded independent random variables}
\author{Jackson Loper\\
Department of Statistics, University of Michigan\\
and \\
Jeffrey Regier \\
Department of Statistics, University of Michigan}
\maketitle

\bigskip
\begin{abstract}
We construct a new tail bound for the sum of independent random variables for situations in which the expected value of the sum is known and each random variable lies within a specified interval, which may be different for each variable.  This new bound can be computed by solving a two-dimensional convex optimization problem.  Simulations demonstrate that the new bound is often substantially tighter than Hoeffding's inequality for cases in which both bounds are applicable.
\end{abstract}

\noindent%
{\it Keywords:}  Hoeffding's inequality; tail bounds; convex optimization
\vfill

\newpage
\spacingset{1.45} % DON'T change the spacing!




\section{Introduction}

The celebrated Hoeffding's inequality \citep[Theorem 2]{hoeffding_probability_1963} has been a cornerstone for probability and statistics over the last 60 years.  This inequality, which we refer to as the ``general Hoeffding inequality,'' concerns the sum of $n$ independent random variables, $(X_1,\ldots, X_n)$.  Each variable $X_i$ lies in some known interval $[a_i,b_i]$, for $i\in \{1,\ldots n\}$.  In this article, without loss of generality, we assume that $a_i=0$ for each $i$.  The general Hoeffding inequality states that the sum of these random variables, $S=\sum_{i=1}^n X_i$, is unlikely to greatly exceed its expected value, $\mu=\mathbb{E}[S]$.  In particular, for any $t\geq 0$,
\begin{equation} \label{eq:hoeffgeneral}
\mathbb{P}(S \geq \mu+t) \leq \exp\left(-2t^2 / \sum_{i=1}^n b_i^2 \right).
\end{equation}
Hoeffding's seminal work also included a less well-known bound, which is applicable only if every bounding interval is of the same length and the expected value $\mu$ is known \citep[Theorem 1]{hoeffding_probability_1963}.  If $b_i=1$ for every $i$, then
\begin{equation} \label{eq:hoeffspecial}
\mathbb{P}(S \geq \mu+t) \leq \left(\frac{\mu}{\mu+t}\right)^{n\mu+nt} \left(\frac{1-\mu}{1-\mu-t}\right)^{n-n\mu-nt}.
\end{equation}
We refer to Equation~\ref{eq:hoeffspecial} as the ``specialized Hoeffding inequality.'' Equation~\ref{eq:hoeffspecial} is tighter than Equation~\ref{eq:hoeffgeneral} in situations where both are applicable.

In this article, we propose a computable bound for the sum of independent random variables for a setting in which the expected value of the sum is known and each random variable is lies in a known fixed interval of arbitrary size.  The size of these intervals is permitted to differ among variables, which makes the direct application of the specialized Hoeffding inequality impossible.  Bounds for this setting are useful for hypothesis testing regarding the sum's expected value based on an observed realization (or, equivalently, for computing confidence intervals on the sum's expected value).  Our key insight is that a worst-case Chernoff--Cram\'er bound related to our setting can be computed exactly by solving a two-dimensional convex optimization problem, yielding a tighter bound than the general Hoeffding inequality.  The result reduces to Equation~\ref{eq:hoeffspecial} in the special case that $b_i=1$ for each $i \in \{1,\ldots,n\}$.  We use simulations to demonstrate that the new bound can be much tighter than the general Hoeffding inequality in situations where both are applicable and the specialized Hoeffding inequality is not.

Note that the generalied Hoeffding bound, the specialized Hoeffding bound, and the bounds we develop here for our setting can all be applied to translated intervals of the form $[a_i, b_i]$ without difficulty: we can always translate each variable by defining $Y_i = X_i - a_i$, which lies in $[0, b_i - a_i]$, and the sum $\sum_i Y_i = \sum_i X_i - \sum_i a_i$ differs from the original sum by a known constant.  The tail bound for $S = \sum_i X_i$ then follows directly from the bound for $\sum_i Y_i$.  This reduction is possible because we are interested in tail bounds of the form $\mathbb{P}(S \geq s)$; the translation preserves the structure of the problem.

\section{Related Work}

We here discuss two related lines of work.

One related line of work considers the case where it is known that $a_i \leq X_i - \mathbb{E}[X_i] \leq b_i$ for every $i$, with both $a_i$ and $b_i$ known per variable.  We refer to this as the ``bounded mean deviation'' (BMD) setting.  Note that although in our setting we may take $a_i=0$ without loss of generality, in the BMD setting the values $a_i$ cannot be simply shifted away.  Instead, the magnitude of $a_i<\leq0$ and $b_i>\geq0$ provide crucial insight into skew for the BMD setting.

Bounds from our setting (in which $a_i \leq X_i \leq b_i$) can be applied to the BMD setting to calculate the probability of a deviation from the mean: apply a method from our setting to $Y_i = X_i - \mathbb{E}[X_i]$ to bound $\mathbb{P}(\sum_{i}^n Y_i > t)$, then translate by $\mu$ to obtain a bound on $\mathbb{P}(\sum_{i}^n X_i > \mu+t)$.  The application of the generalized Hoeffding inequality to this BMD setting has been shown to be unnecessarily loose.  For example, \citet{hertz2020improved} notes that if $a_i\neq -b_i$ then the asymmetry can be used to produce a superior bound.  On the other hand, when $a_i=-b_i$, \citet[Cor 1.4]{bentkus2004hoeffding} provides a different bound that is often superior to the generalized Hoeffding inequality.

A bound derived from a BMD assumption can be applied to our setting, but the bounds are necessarily quite loose.  To see why, consider a simple version of our setting in which $n=100$ and $X_i \in [-1/2,1/2]$ for each $i$.  This implies the BMD assumption that $1 \leq X_i - \mathbb{E}[X_i] \leq 1$.  However, this BMD assumption is much weaker than our original assumption, and it yields accordingly poor bounds.  In particular, the implied BMD assumption is compatible a case where $X_i$ carries the uniform distribution on $\{-1,1\}$, so that, for example, $\log \mathbb{P}(S>20) \approx -4.03985$.  Consequently, any upper bound based solely on this BMD assumption must be at least $-4.04$.  On the other hand Hoeffding's generalized inequality based on the original assumption yields $\mathbb{P}(S>\mu+20) \leq -8$.  Note that this discrepancy between $-4.04$ and $-8$ is only possible because the uniform distribution producing the $-4.04$ value actually \emph{violates} our original assumption.  In general, the large-deviations exponent for bounds derived from the BMD assumption alone will tend to underperform the large-deviations exponent of bounds derived from the original assumption by a factor of two.

Another line of work considers the case where each $X_i \in [0,1]$.  We refer to this as the ``constant width'' (CW) setting.  The CW setting is a special case of our setting.  Bounds from the CW setting can be applied to our setting by applying a single shift and rescaling across all variables such that they lie in the unit interval.  However, this yields necessarily loose bounds, as it throws away information about the per-variable bounds.  The original specialized Hoeffding inequality targeted this setting, but bounds for this setting have been improved in the meanwhile, for example by \citet[Thm 1.2]{bentkus2004hoeffding}.  In simulations we compare the performance of our method with this method of Bentkus in the CW case (note that our method reduces to Hoeffding's specialized inequality in the CW case).  There has also been some work investigating how additional assumptions can enable tighter bounds in the CW case.  For example, \citet{from2013refinement} show a refinement if $\mu_i=\mathbb{E}[X_i]$ is known for each $i$, and \citet{from2024refinement} show a refinement if $\beta = \sum_i^n (\mu_i-\mu)^2/n$ is known.

\section{Notation and preliminaries}

\label{sec:notation}

Consider a set of probability distributions on independent bounded random variables.  For any scalar $c\geq 0$, let $\mathcal{M}_c$ denote the space of probability distributions with support on $[0,c]$.  For any vector of interval lengths $\boldsymbol{b}\in ({\mathbb{R}^+})^n$ and fixed value $\mu$, let
\begin{equation} \label{eq:mdef}
\mathcal{M}_{\boldsymbol{b},\mu} = \left\{\boldsymbol{p} \in \prod_{i=1}^n  \mathcal{M}_{b_i}:\ \sum_{i=1}^n \mathbb{E}_{X_i \sim p_i}[X_i] \leq \mu\right\}.
\end{equation}
In other words, a random vector $\boldsymbol{X}=(X_1,\ldots X_n)$  drawn according to a member of $\mathcal{M}_{\boldsymbol{b},\mu}$ satisfies three conditions: (i) each variable is independent, (ii) $X_i \in [0,b_i]$ for each $i$, and (iii) the expected sum over the vector is at most $\mu$.

\citet[Theorem 2]{hoeffding_probability_1963} provides a tail bound for the variable $S=\sum_{i=1}^n X_i$ when $\boldsymbol{X}$ is drawn according to $\mathcal{M}_{\boldsymbol{b},\mu}$, for fixed $\boldsymbol{b}$ and $\mu$.  To construct this bound, \citet{hoeffding_probability_1963} first considers a tail bound for a particular $\boldsymbol{p} \in \mathcal{M}_{\boldsymbol{b},\mu}$.
Let
\[
\varphi(\boldsymbol{p},s) = \inf_{t\geq 0} \left(\sum_{i=1}^n \log \mathbb{E}_{X_i \sim p_i}[\exp(t X_i)] - ts\right).
\]
If $\boldsymbol{X}$ is drawn according to any member of $\mathcal{M}_{\boldsymbol{b},\mu}$, the Chernoff--Cram\'er bound gives that $\log \mathbb{P} \left (S \geq s\right)$ is at most $\sup_{\boldsymbol{p} \in \mathcal{M}_{\boldsymbol{b},\mu}} \varphi(\boldsymbol{p},s)$.
To produce a bound that is universal for all members of $\mathcal{M}_{\boldsymbol{b},\mu}$, \citet{hoeffding_probability_1963} therefore considers the quantity
\begin{equation} \label{eq:varphistar}
\varphi^*_{\boldsymbol{b},\mu}(s) = \sup_{\boldsymbol{p} \in \mathcal{M}_{\boldsymbol{b},\mu}} \varphi(\boldsymbol{p},s) = \sup_{\boldsymbol{p} \in \mathcal{M}_{\boldsymbol{b},\mu}} \inf_{t\geq 0} \left(\sum_{i=1}^n \log \mathbb{E}_{X_i \sim p_i}[\exp(t X_i)] - ts\right).
\end{equation}
This leads to a tail bound for $S$ in terms of $\varphi^*$:
\begin{equation} \label{eq:chernoff}
\mathbb{P} \left (S \geq s\right) \leq \exp(\varphi^*_{\boldsymbol{b},\mu}(s)).
\end{equation}
Equation~\ref{eq:chernoff} is the worst-case Chernoff--Cram\'er bound for a sum of independent random variables $(X_1,\ldots X_n)$ with each $X_i \in [0,b_i]$ and $\sum_{i=1}^n \mathbb{E}[X_i]=\mu$.

To compute the right-hand side of Equation~\ref{eq:chernoff}, \citet{hoeffding_probability_1963} considers two different readily computable expressions for $\varphi^*$.  In one theorem, a formula for $\varphi^*_{\boldsymbol{b},\mu}$ is provided for the special case in which $b_1=\cdots=b_n=1$, leading to the specialized Hoeffding inequality, Equation~\ref{eq:hoeffspecial}.  In another theorem, an \emph{upper bound} is produced for $\varphi^*_{\boldsymbol{b},\mu}$ in the general case of arbitrary $\boldsymbol{b}$, namely,
$
\varphi^*_{\boldsymbol{b},\mu}(s)\leq -2(s-\mu)^2/\sum_{i=1}^n b_i^2,
$
leading to the general Hoeffding inequality.

Our contribution in the next section is to show how $\varphi^*$ can be \textit{exactly} computed for arbitrary interval lengths $\boldsymbol{b}$.
By substituting this exact value of $\varphi^*$ into Equation~\ref{eq:chernoff}, an improved tail bound for the sum of bounded independent random variables follows.

\section{Optimizing the worst-case Chernoff--Cram\'er bound}

We now show that $\varphi^*$ (cf. Equation~\ref{eq:varphistar}) can be computed by solving a two-dimensional convex optimization problem.  We begin by expressing $\varphi^*$ as the solution to a high-dimensional convex optimization problem and then demonstrate how to reduce this problem to a more tractable two-dimensional one.

Equation~\ref{eq:varphistar} defines $\varphi^*$ as the solution to a maximin problem that involves maximizing over distributions and minimizing over values of $t>0$.  The inner one-dimensional minimization over $t$ can be commuted with the outer infinite-dimensional maximization over $\boldsymbol{p}\in\mathcal{M}_{\boldsymbol{b},\mu}$.  Once the minimization and maximization operators have been commuted, the inner problem reduces to a finite-dimensional optimization problem.  This leads to a new expression for $\varphi^*$, as shown below.

\begin{theorem} \label{thm:tightchernoff} Fix $\boldsymbol{b}=(b_1,\ldots b_n)$, $\mu \in [0,\sum_{i=1}^n b_i]$, and $s \in [0,\sum_{i=1}^n b_i]$.  Let
\begin{align*}
T=\left\{\tau\in \prod_{i=1}^n [0,b_i] : \sum_{i=1}^n \tau_i = \mu\right\}.
\end{align*}
Then, $\varphi^*_{\boldsymbol{b},\mu}(s)$ from Equation~\ref{eq:varphistar} can be expressed as
\begin{gather*}
\varphi^*_{\boldsymbol{b},\mu}(s)
= \min_{t\geq 0} \left(\max_{\tau \in T} \sum_{i=1}^n \log\left(1+\tau_i (\exp(b_i t)-1)/b_i \right) - ts\right).
\end{gather*}
\end{theorem}

The high-dimensional minimax problem from Theorem~\ref{thm:tightchernoff} can be reduced to a two-dimensional convex optimization problem.

\begin{theorem} \label{thm:lowd} Fix $\boldsymbol{b}=(b_1,\ldots b_n)$, $\mu \in [0,\sum_{i=1}^n b_i]$, and $s \in [0,\sum_{i=1}^n b_i]$.  Let $\xi(b,t) = (\exp(bt)-1)/b$,
\[
\tau^*_i(t,\lambda) = \min\left(\max\left(\frac{\xi(b_i,t)-\lambda}{\xi(b_i,t)\lambda},0\right),b_i\right),
\]
and
\[
g(t,\lambda;s) = \sum_i \log\left(1+\xi(b_i,t)\tau_i^*(t,\lambda)\right) + \lambda\left(\mu - \sum_i \tau_i^*(t,\lambda)\right)- ts.
\]
Then, $g$ is convex and $\varphi^*_{\boldsymbol{b},\mu}(s)$ from Equation~\ref{eq:varphistar} can be expressed as
\[
\varphi^*_{\boldsymbol{b},\mu}(s) = \min_{t,\lambda \geq 0} g(t,\lambda;s).
\]
Moreover, the mapping $t \mapsto \min_\lambda g(t,\lambda;s)$ is convex and the mapping $\lambda \mapsto g(t,\lambda)$ is convex for each $t \geq 0$.
\end{theorem}

This theorem suggests two methods for computing $\varphi^*$.  One method involves directly minimizing the two-dimensional convex function $t,\lambda \mapsto g(t,\lambda;s)$ and the other involves minimizing the one-dimensional convex function $t \mapsto \min_\lambda g(t,\lambda;s)$. Because each evaluation of the objective function of the latter requires solving a different one-dimensional convex optimization problem, in our simulations, we use the first method.

\section{Simulations}

\label{sec:simulations}

\begin{figure}
\begin{center}
\includegraphics[width=5in]{comparisons/dirichlet_100.png}
  \caption{
  Comparison of bounds for heterogeneous intervals.  We consider three choices of fixed expected value, $\mu \in \{0.8, 0.9, 0.95\}$.  We compare the general Hoeffding inequality \citep[Theorem 2]{hoeffding_probability_1963} (dashed lines), our new Tight Chernoff bounds (solid lines, calculated using Theorem~\ref{thm:lowd}), Bentkus Theorem 1.2 \citep{bentkus2004hoeffding} (dotted lines), and Bentkus Corollary 1.4 (×~markers).  Since Bentkus Theorem 1.2 requires homogeneous intervals, we apply it by rescaling all variables to lie within $[0, \max_i b_i]$.  Bentkus Corollary 1.4 uses RMS scale but requires centering information.  Our Tight Chernoff bound substantially outperforms both the general Hoeffding inequality and the Bentkus bounds in the heterogeneous case.
  }\label{fig:tailboundcomparisons}
\end{center}
\end{figure}

We use simulations to assess the degree of improvement over the general Hoeffding inequality.  First, we sample a vector of interval lengths, $\boldsymbol{b} \in \mathbb{R}^{100}$, uniformly at random from the set of lengths such that $B=\sum_{i=1}^{100} b_i=1$.  We define $S$ to be the sum of 100 independent random variables, $(X_1,\ldots,X_{100})$, where each $X_i$ is constrained to the interval $[0,b_i]$ for $i\in \{1,\ldots,100\}$.  The value of $n$ affects the tightness of all bounds.  However, the overall scaling of the $b_i$ values does not matter: what is relevant is the relative sizes of the intervals, so we normalize the intervals so that their maximum possible sum is 1.

For any fixed value of $\mu$ and any value $s \geq \mu$, our new results provide an upper bound on $\mathbb{P}(S>s)$, assuming $\mathbb{E}[S]=\mu$.  The general Hoeffding inequality can also be used to produce a bound on $\mathbb{P}(S>s)$, offering a natural point of comparison.  We also compare with the two Bentkus bounds \citep{bentkus2004hoeffding}: Theorem 1.2 requires all bounding intervals to be of the same size, while Corollary 1.4 uses RMS scale but requires centering information.  To apply Theorem 1.2 to heterogeneous intervals, we rescale all variables so that they lie within $[0, \max_i b_i]$.  For Corollary 1.4, we apply it using only the overall centering $S - \mu$.  These bounds on $\mathbb{P}(S>s)$ have practical implications, such as facilitating hypothesis testing for $\mathbb{E}[S]=\mu$ given an observed realization of $S$.

Figure~\ref{fig:tailboundcomparisons} compares our new Tight Chernoff bounds with the general Hoeffding inequality and both Bentkus bounds for $\mu \in \{0.8, 0.9, 0.95\}$, showing how they vary as a function of $s$. In many cases, our bounds are orders of magnitude tighter than those derived from the general Hoeffding inequality, demonstrating their advantage in tail estimation.  Neither Bentkus bound performs as well as our method in the heterogeneous case: Theorem 1.2 effectively treats all intervals as having the maximum width, while Corollary 1.4 has a larger multiplicative constant and does not fully exploit the heterogeneity.

\begin{figure}
\begin{center}
\includegraphics[width=5in]{comparisons/homogen_100.png}
  \caption{
  Comparison of bounds for homogeneous intervals.  When all variables share the same bounding interval $[0,1]$, the specialized Hoeffding inequality (Equation~\ref{eq:hoeffspecial}) applies.  We compare the general Hoeffding inequality, the specialized Hoeffding inequality, our new Tight Chernoff bounds, Bentkus Theorem 1.2, and Bentkus Corollary 1.4.  Our Tight Chernoff bounds match the specialized Hoeffding inequality, as expected from Theorem~\ref{thm:lowd}.  In this homogeneous setting, the specialized Hoeffding inequality, the Bentkus Theorem 1.2 bound, and our Tight Chernoff bounds all provide much tighter bounds than the general Hoeffding inequality.  Bentkus Corollary 1.4 is less tight here due to its larger multiplicative constant ($2e^3/9 \approx 4.46$ vs.\ $e \approx 2.72$ for Theorem 1.2).
  }\label{fig:homogencomparisons}
\end{center}
\end{figure}

To verify that our bounds perform well even when the Bentkus bounds are directly applicable, we also consider the case of homogeneous intervals where all $b_i = 1$.  Figure~\ref{fig:homogencomparisons} shows this comparison.  In this setting, our new Tight Chernoff bounds coincide with the specialized Hoeffding inequality (Equation~\ref{eq:hoeffspecial}), as guaranteed by Theorem~\ref{thm:lowd}.  Both our bounds and Bentkus Theorem 1.2 are much tighter than the general Hoeffding inequality, demonstrating that all refined methods provide substantial improvements when applicable.  Bentkus Corollary 1.4 is less tight in this homogeneous case because its larger prefactor ($2e^3/9 \approx 4.46$) outweighs any benefit from using RMS scale (which equals the max scale when intervals are equal).  The key advantage of our Tight Chernoff approach is that it extends naturally to heterogeneous intervals without requiring any rescaling or centering information beyond the overall mean.


\begin{appendix}
\section*{Proofs}\label{appn}

To prove Theorem~\ref{thm:tightchernoff}, it is convenient to view $\mathcal{M}_{\boldsymbol{b},\mu}$ from Equation~\ref{eq:mdef} as a topological vector space.  For any scalar $b\geq 0$, let $\mathcal{M}_b$ denote the space of probability measures with support on $[0,b]$.  Endowing $\mathcal{M}_b$ with the weak-$*$ topology, we view $\mathcal{M}_b$ as a topological vector space.  Note that $\mathcal{M}_b$ is compact and convex.  For any values $\boldsymbol{b}\in ({\mathbb{R}^+})^n$ and fixed expected value $\mu$, let
\[
\mathcal{M}_{\boldsymbol{b},\mu} = \left\{\boldsymbol{p} \in \prod_i \mathcal{M}_{b_i}:\ \sum_i\mathbb{E}_{X_i \sim p_i}[X_i] \leq \mu\right\}.
\]
Here
$\prod_i \mathcal{M}_{b_i}$
denotes a direct sum of a topological vector spaces.  For example, if $\boldsymbol{p} \in \mathcal{M}_{\boldsymbol{b},\mu}$ then $\boldsymbol{p}=(p_1,\ldots, p_n)$ where each $p_i \in \mathcal{M}_{b_i}$.  Recall that addition is computed in the following manner for such direct sum spaces: for any $\boldsymbol{p},\boldsymbol{q} \in \mathcal{M}_{\boldsymbol{b},\mu}$, $(\boldsymbol{p}+\boldsymbol{q})=((p_1+q_1),\ldots, (p_n+q_n))$.  This seemingly excessive formality is necessary to clarify that $\mathcal{M}_{\boldsymbol{b},\mu}$ is, in fact, convex.

\begin{theorem*}[Restatement of Theorem 3.1] Fix $\boldsymbol{b}=(b_1,\ldots b_n)$, $\mu \in [0,\sum_{i=1}^n b_i]$, and $s \in [0,\sum_{i=1}^n b_i]$.  Let $T=\{\tau\in \prod_{i=1}^n [0,b_i]: \sum_{i=1}^n \tau_i = \mu\}$.   Then $\varphi^*_{\boldsymbol{b},\mu}(s)$ from Equation~\ref{eq:varphistar} can be expressed as
\begin{gather*}
\varphi^*_{\boldsymbol{b},\mu}(s)
= \min_{t\geq 0} \left(\max_{\tau \in T} \sum_{i=1}^n \log\left(1+\tau_i (\exp(b_i t)-1)/b_i \right) - ts\right).
\end{gather*}
\end{theorem*}

\begin{proof}

Fix $s$. Let $\xi(b,t)=(\exp(bt)-1)/b$.  Let
\[
f(t,\boldsymbol{p})=\sum_{i}\log \mathbb{E}_{B_{i}\sim p_{i}}[\exp(tB_{i})]-ts.
\]
Our object of interest is may be given as $\min_{t}\max_{\boldsymbol{p}}f(t,\boldsymbol{p})$.
Our first step is to reverse the order of the minimization and maximization
by applying Sion's minimax theorem \citep{komiya1988elementary}. We observe the following.
\begin{itemize}
\item $f$ is convex with respect to $t$, as it is the sum of cumulant
generating functions.
\item $f$ is concave with respect to $\boldsymbol{p}$, due to Jensen's
inequality.
\item $f$ is continuous with respect both $\boldsymbol{p}$ and $t$.
\item $\mathcal{M}_{\boldsymbol{b},\mu}$ is both convex and compact.
\item $[0,\infty)$ is convex (though not compact).
\end{itemize}
Sion's minimax theorem thus shows that
\[
\inf_{t\geq0}\sup_{\boldsymbol{p}\in\mathcal{M}_{\boldsymbol{b},\mu}}f(t,\boldsymbol{p})=\sup_{\boldsymbol{p}\in\mathcal{M}_{\boldsymbol{b},\mu}}\inf_{t\geq0}f(t,\boldsymbol{p}).
\]
For any fixed $t$, the maximization of $f$ may be reduced to a finite-dimensional
problem. First, we rewrite the maximization problem by introducing
auxiliary variables, as follows:
\begin{align*}
\max_{\boldsymbol{p},\tau}\qquad & f(t,\boldsymbol{p})\\
\mathrm{s.t.}\qquad & \tau_{i}\in[0,b_i]\\
 & p_{i}\in\mathcal{M}_{b_i}\\
 & \tau_{i}=\sum_{i}\mathbb{E}_{B_{i}\sim p_{i}}\left[B_{i}\right]\\
 & \sum_{i}\tau_{i} \leq \mu
\end{align*}
For any fixed $\tau$, Hoeffding \citep[proof of Theorem 1]{hoeffding_probability_1963} showed that the cumulant generating
functions can be maximized by setting each $p_{i}$ to be a scaled
Bernoulli random variable, namely
\[
p_{i}=\delta_{0}\left(1-\frac{\tau_{i}}{b_i}\right)+\delta_{b_i}\frac{\tau_{i}}{b_i}
\]
where $\delta_{x}$ represents the point mass at $x$. Under this
distribution,
\[
\log \mathbb{E}_{p_{i}}\left[\exp\left(tB_{i}\right)\right]=\log\left(1+\xi(b_i,t)\tau_{i}\right).
\]
Summing over $i$, we obtain that
\begin{gather*}
\varphi^*_{\boldsymbol{b},\mu}(s)
= \min_{t\geq 0} \left(\max_{\tau\in \prod_{i=1}^n [0,b_i]: \sum_{i=1}^n \tau_i \leq \mu} \sum_{i=1}^n \log\left(1+\tau_i (\exp(b_i t)-1)/b_i \right) - ts\right).
\end{gather*}
This is very nearly the statement we seek to prove.  The only remaining impediment is that the maximization over $f$ that we
have defined permits $\sum_i\tau_i \leq \mu$, whereas we claim in the theorem statement that it is sufficient to only consider
cases in which $\sum_i\tau_i = \mu$.

To see that this is true, note that the objective is coordinatewise nondecreasing in $\tau$ (strictly increasing if $t>0$).  Hence, if a maximizer
$\tau^\star$ under $\sum_i\tau_i\le \mu$ had slack, i.e.\ $\sum_i\tau_i^\star<\mu$, then since
$\mu\le\sum_i b_i$ there exists some $j$ with $\tau_j^\star<b_j$; increasing $\tau_j^\star$ by
$\varepsilon=\min\{b_j-\tau_j^\star,\ \mu-\sum_i\tau_i^\star\}>0$ preserves feasibility and
(in case $t>0$) strictly increases $F$, a contradiction.  Thus the maximum is attained at
$\sum_i\tau_i=\mu$ (and when $t=0$ the objective is constant, so the value is unchanged regardless of $\tau$).

\end{proof}

\begin{theorem*}[Restatement of Theorem 3.2] Fix $\boldsymbol{b}=(b_1,\ldots b_n)$, $\mu \in [0,\sum_{i=1}^n b_i]$, and $s \in [0,\sum_{i=1}^n b_i]$.  Let $\xi(b,t) = (\exp(bt)-1)/b$,
\[
\tau^*_i(t,\lambda) = \min\left(\max\left(\frac{\xi(b_i,t)-\lambda}{\xi(b_i,t)\lambda},0\right),b_i\right),
\]
and
\[
g(t,\lambda;s) = \sum_i \log\left(1+\xi(b_i,t)\tau_i^*(t,\lambda)\right) + \lambda\left(\mu - \sum_i \tau_i^*(t,\lambda)\right)- ts.
\]
Then, $g$ is convex and $\varphi^*_{\boldsymbol{b},\mu}(s)$ from Equation~\ref{eq:varphistar} can be expressed as
\[
\varphi^*_{\boldsymbol{b},\mu}(s) = \min_{t,\lambda \geq 0} g(t,\lambda;s).
\]
Moreover, the mapping $t \mapsto \min_\lambda g(t,\lambda;s)$ is convex and the mapping $\lambda \mapsto g(t,\lambda)$ is convex for each $t \geq 0$.
\end{theorem*}

\begin{proof}
Let $T=\{\tau\in \prod_{i=1}^n [0,b_i]: \sum_{i=1}^n \tau_i = \mu\}$.  Recall from Theorem~\ref{thm:tightchernoff} that
\begin{gather*}
\varphi^*_{\boldsymbol{b},\mu}(s)
= \min_{t\geq 0} \left(\max_{\tau \in T} \sum_{i=1}^n \log\left(1+\tau_i (\exp(b_i t)-1)/b_i \right) - ts\right).
\end{gather*}
For any fixed $t\geq0$, we begin by rewriting the inner maximization problem
in terms of
\[
f(\tau;t)=\begin{cases}
\sum_{i}\log\left(1+\xi(b_i,t)\tau_{i}\right)-ts & \mathrm{if}\ \tau\in\prod_{i}[0,b_i]\\
-\infty & \mathrm{else}.
\end{cases}
\]
The inner maximization problem is thus equivalent to
\[
\max_{\tau:\ \sum_{i}\tau_{i}=\mu}f(\tau;t).
\]
The associated Lagrangian function is given by
\[
\mathcal{L}(\tau,\lambda;t)=\sum_{i}\log\left(1+\xi(b_i,t)\tau_{i}\right)+\lambda\left(\mu-\sum_{i}\tau_{i}\right)-ts.
\]
For any fixed $t\geq0$, observe that $\xi(b_i,t)$ is positive
and thus $\tau\mapsto\mathcal{L}\left(\tau,\lambda;t\right)$ is concave.
The argument maximizing $\tau\mapsto\mathcal{L}(\tau,\lambda;t)$
is given by $\tau_{i}^{*}(t,\lambda)$. Thus $g(t,\lambda)=\max_{\tau}\mathcal{L}(\tau,\lambda;t)$,
the Lagrangian dual for a convex optimization problem with convex
constraints. There is at least one feasible point, namely $\tau_{i}=b_i\mu/\sum_{j}b_{j}$,
so strong duality holds and
\[
\max_{\substack{\tau\in\prod_{i}[0,b_i]\\
\sum_{i}\tau_{i}=\mu
}
}\sum_{i}\log\left(1+\xi(b_i,t)\tau_{i}\right)-ts=\min_{\lambda}g(t,\lambda)
\]
as desired.

We now demonstrate that $g$ is convex. First, observe that $\mathcal{L}$
is affine in $\lambda$ and convex in $t$. Thus $g$ is a pointwise
maximum of a family of convex functions: it is convex.

Finally, we demonstrate that $t\mapsto\min_{\lambda}g(t,\lambda)$
is convex. Observe that, for any fixed $\tau_{i}\geq0$, the mapping
\[
t\mapsto\sum_{i}\log\left(1+\xi(b_i,t)\tau_{i}\right)-ts
\]
is convex in $t$. Thus $\min_{\lambda}g(t,\lambda)$ is also the
pointwise maximum of a family of convex functions: it is also convex.
\end{proof}

\section*{Sensitivity to Sample Size}\label{app:sensitivity}

In this appendix, we present additional simulation results to demonstrate how the bounds vary with different sample sizes.  The main text presents results for $n=100$; here we show analogous results for $n=50$ and $n=200$.

\begin{figure}[h]
\begin{center}
\includegraphics[width=5in]{comparisons/dirichlet_50.png}
  \caption{Comparison of bounds for heterogeneous intervals with $n=50$ variables.}\label{fig:dirichlet50}
\end{center}
\end{figure}

\begin{figure}[h]
\begin{center}
\includegraphics[width=5in]{comparisons/homogen_50.png}
  \caption{Comparison of bounds for homogeneous intervals with $n=50$ variables.}\label{fig:homogen50}
\end{center}
\end{figure}

\begin{figure}[h]
\begin{center}
\includegraphics[width=5in]{comparisons/dirichlet_200.png}
  \caption{Comparison of bounds for heterogeneous intervals with $n=200$ variables.}\label{fig:dirichlet200}
\end{center}
\end{figure}

\begin{figure}[h]
\begin{center}
\includegraphics[width=5in]{comparisons/homogen_200.png}
  \caption{Comparison of bounds for homogeneous intervals with $n=200$ variables.}\label{fig:homogen200}
\end{center}
\end{figure}

As the sample size increases, all bounds become tighter, but the relative advantage of the Tight Chernoff bound over the general Hoeffding inequality remains substantial across all values of $n$ considered.

\end{appendix}



\bigskip
\begin{center}
{\large\bf SUPPLEMENTARY MATERIAL}
\end{center}

\begin{description}

\item[Demonstrating the new bound:] Code for solving the two-dimensional convex optimization problem involved in the new bound, together with demonstrations of its use.  (Jupyter notebook)

\end{description}

\bibliographystyle{agsm}
\bibliography{bibliography}       % Bibliography file (usually '*.bib')


\end{document}
